\chapter{Часто задаваемые вопросы (FAQ)}
\label{ch:faq}
\index{FAQ}

\begin{description}
  \item[Зачем Pipenv и \texttt{make install}?]\hfill\\
  Зависимости Python должны устанавливаться только в изолированное окружение (см.\ главу~\ref{ch:infra}), а не в системный интерпретатор. \texttt{make install} создаёт виртуальное окружение и ставит пакеты из \texttt{Pipfile}.

  \item[Чем отличаются \texttt{src\_starting\_point} и \texttt{src\_solution}?]\hfill\\
  \texttt{src\_starting\_point} --- заготовка АБУ (отправная точка; не претендует на оптимальность архитектуры). Решение участника размещается в \texttt{src\_solution/} (или в согласованном с жюри каталоге).

  \item[На диаграммах в PDF пустые рамки вместо рисунков. Что делать?]\hfill\\
  Сгенерируйте PNG из PlantUML в корне репозитория командой \texttt{make diagrams}, затем пересоберите PDF из каталога \texttt{docs/bundle}.

  \item[Как сдать работу?]\hfill\\
  Обычно требуется ссылка на \textbf{приватный} git-репозиторий и добавление пользователя организаторов с правами разработчика --- уточняйте в объявлении жюри.

  \item[Где полный перечень критериев C01--C25?]\hfill\\
  В репозитории прототипа: регламент и скрипт оценки; в настоящем документе --- глава~\ref{ch:evaluation}.

  \item[Обязателен ли Docker?]\hfill\\
  Нет: основные проверки выполняются через \texttt{make tests-all} локально. Docker нужен для сценариев интеграции с поднятыми сервисами ЦР/Регулятора в контейнерах.

  \item[Что такое \texttt{certificate\_id}?]\hfill\\
  В прототипе это идентификатор, связанный с результатом сертификации пакета (см.\ главу~\ref{ch:certification}); используйте значение из ответа Регулятора при регистрации АБУ в ЦР.
\end{description}
